\section{$q$-展开的推导}
\label{sec:eisenstein-q-expansion}
% ============================================================================

有了 Poisson 求和对 $S_c(\tau)=\sum_{d\in\Z}(c\tau+d)^{-k}$ 的计算，最难的分析步骤其实已经完成了。
现在我们要做的，是把这些局部碎片重新拼起来，得到 Eisenstein 级数整体的 $q$-展开。
不过在正式拼装之前，先别急着算；更重要的是先想清楚：\textbf{为什么我们如此在意 $q$-展开？}

对模形式来说，$q$-展开就像它的 DNA。
一旦把 $f(\tau)$ 写成 $\sum a_n q^n$，许多原本抽象的事情都会立刻变得具体：
做加法、乘法时，只是在处理幂级数；比较两个模形式是否相等时，只要比较它们的系数；研究 Hecke 算子时，作用也会直接落在这些系数上。
更重要的是，真正有算术味道的信息几乎都藏在 $a_n$ 里面——例如 Eisenstein 级数的系数会变成 $\sigma_{k-1}(n)$，把整数的因子结构编码进去。

所以这一节的目标不只是写出一个好看的公式，而是把 Eisenstein 级数从"定义出来的解析函数"，转变成"可以计算、可以比较、可以承载算术信息的对象"。
而我们即将得到的公式
\[
  E_k(\tau)=1-\frac{2k}{B_k}\sum_{n=1}^{\infty}\sigma_{k-1}(n)q^n
\]
也会顺便解释一个很自然的问题：为什么系数偏偏是除数和？
原因就在上一节的 Fourier 变换里：频率 $m$ 给出 $m^{k-1}$，再把所有满足 $mc=n$ 的配对收集起来，就恰好形成
$\sum_{m\mid n} m^{k-1}=\sigma_{k-1}(n)$。
这正是分析变换如何产出算术函数的一个漂亮样板。

% figures/ch04-q-expansion-geometry.tex
% q 展开的几何意义：τ → q = e^{2πiτ} 把竖带映到穿孔圆盘
\begin{figure}[ht]
\centering
\begin{tikzpicture}[>=Stealth, font=\small, scale=1.1]

% === 左侧：上半平面的竖带 ===
\begin{scope}[xshift=-4.2cm]
  % 阴影填充带
  \fill[blue!8]
    (-1, 0.3) rectangle (1, 3.5);
  % 边界
  \draw[thick, blue!60!black] (-1, 0.3) -- (-1, 3.5);
  \draw[thick, blue!60!black] (1, 0.3) -- (1, 3.5);
  \draw[dashed, gray] (-1.5, 0.3) -- (1.5, 0.3)
    node[right, font=\footnotesize] {$\im(\tau) = Y$};
  % 坐标轴
  \draw[->] (-1.6, 0) -- (1.6, 0) node[right] {$\Re(\tau)$};
  \draw[->] (0, 0) -- (0, 3.8) node[above] {$\im(\tau)$};
  % 刻度
  \node[below] at (-1, 0) {$-\tfrac{1}{2}$};
  \node[below] at (1, 0) {$\tfrac{1}{2}$};
  % 示例点
  \fill[red!70!black] (0.3, 1.5) circle (2pt);
  \node[right, red!70!black, font=\footnotesize] at (0.33, 1.5) {$\tau$};
  % 顶部箭头（∞方向）
  \draw[->, green!50!black, line width=1.2pt] (0, 3.2) -- (0, 3.7);
  \node[green!50!black, right, font=\footnotesize] at (0.05, 3.5) {$\infty$};
  % 标签
  \node[blue!60!black] at (0, 0.9) {$\HH$（竖带）};
\end{scope}

% === 中间箭头 ===
\draw[->, thick, purple!70!black, line width=1.5pt]
  (-1.8, 1.8) -- (-0.3, 1.8)
  node[above, midway, font=\footnotesize, purple!70!black]
  {$q = e^{2\pi i\tau}$};

% === 右侧：穿孔圆盘 ===
\begin{scope}[xshift=1.8cm]
  % 圆盘
  \fill[orange!10] (0,0) circle (2cm);
  \fill[white] (0,0) circle (0.25cm);  % 穿孔
  \draw[thick, orange!70!black] (0,0) circle (2cm);
  \draw[thick, dashed, orange!40!black] (0,0) circle (0.25cm);
  % r=|q|=e^{-2πY}的圆
  \draw[dashed, blue!50!black] (0,0) circle (1.2cm);
  \node[blue!50!black, font=\footnotesize] at (1.05, 0.65)
    {$|q|=e^{-2\pi Y}$};
  % 示例点 q
  \fill[red!70!black] (0.6, 0.9) circle (2pt);
  \node[right, red!70!black, font=\footnotesize] at (0.63, 0.9) {$q(\tau)$};
  % 穿孔标签
  \node[font=\footnotesize, gray] at (0, -0.45) {$q=0$};
  \draw[->, gray, shorten >=3pt] (0.0, -0.42) -- (0.0, -0.23);
  % 外圆标签
  \node[font=\footnotesize, orange!70!black] at (1.6, -1.6)
    {$|q|=1$};
  \node[font=\footnotesize] at (0, -2.5) {穿孔圆盘 $\mathbb{D}^*$};
\end{scope}

% === 对应关系注记 ===
\node[font=\footnotesize, align=left, text=gray!70!black] at (0, -1.5)
  {$\im(\tau)\to+\infty$ $\Leftrightarrow$ $q\to 0$\quad；\quad
   $\tau \equiv \tau+1$ $\Leftrightarrow$ $q$ 不变（周期性）};

\end{tikzpicture}
\caption{$q$-展开的几何意义。映射 $q = e^{2\pi i\tau}$ 把上半平面中的竖带
$\{|\Re(\tau)| \leq \frac{1}{2},\, \im(\tau) > Y\}$（左）
双全纯地映到穿孔圆盘 $\{0 < |q| < e^{-2\pi Y}\}$（右）。
$\tau \to \infty$ 对应 $q \to 0$，模形式在尖点全纯等价于在 $q=0$ 处可去奇点。}
\label{fig:q-expansion-geometry}
\end{figure}


\begin{theorem}[$G_k$ 的 $q$-展开]
\label{thm:Gk-q-expansion}
对偶整数 $k \geq 4$，令 $q = e^{2\pi i\tau}$（$\tau \in \HH$），则
\[
  G_k(\tau) = 2\zeta(k)
  + \frac{2(2\pi i)^k}{(k-1)!} \sum_{n=1}^{\infty} \sigma_{k-1}(n)\, q^n,
\]
其中 $\sigma_{k-1}(n) = \sum_{d \mid n} d^{k-1}$ 是 $n$ 的因子的 $(k-1)$ 次幂之和。
\end{theorem}

\begin{proof}
由\cref{ex:Gk-first-terms}，
\[
  G_k(\tau) = 2\zeta(k) + 2\sum_{c=1}^{\infty} S_c(\tau),
  \qquad S_c(\tau) = \sum_{d \in \Z} (c\tau + d)^{-k}.
\]
由\cref{ex:poisson-for-eisenstein}（Poisson 求和对 $S_c$），
\[
  S_c(\tau) = \frac{(-2\pi i)^k}{(k-1)!} \sum_{m=1}^{\infty} m^{k-1} e^{2\pi i mc\tau}
  = \frac{(2\pi i)^k}{(k-1)!} \sum_{m=1}^{\infty} m^{k-1} q^{mc}.
\]
（用了 $(-1)^k = 1$ 因为 $k$ 为偶数。）

代入：
\[
  G_k(\tau)
  = 2\zeta(k) + \frac{2(2\pi i)^k}{(k-1)!}
    \sum_{c=1}^{\infty} \sum_{m=1}^{\infty} m^{k-1} q^{mc}.
\]
交换求和顺序，令 $n = mc$（固定 $n$，$m$ 跑遍 $n$ 的正因子，$c = n/m$）：
\begin{align}
  \sum_{c=1}^{\infty} \sum_{m=1}^{\infty} m^{k-1} q^{mc}
  &= \sum_{n=1}^{\infty} \left(\sum_{m \mid n} m^{k-1}\right) q^n
  = \sum_{n=1}^{\infty} \sigma_{k-1}(n)\, q^n.
\end{align}
这就完成了证明。
\end{proof}

\textbf{从 $G_k$ 到 $E_k$。}
用\cref{prop:bernoulli-zeta}，$2\zeta(k) = -\frac{(2\pi i)^k B_k}{k!}$，故
\[
  E_k(\tau)
  = \frac{G_k(\tau)}{2\zeta(k)}
  = 1 - \frac{2k}{B_k} \sum_{n=1}^{\infty} \sigma_{k-1}(n)\, q^n.
\]
（验证：$\frac{2(2\pi i)^k/(k-1)!}{2\zeta(k)} = \frac{2(2\pi i)^k/(k-1)!}{-(2\pi i)^k B_k/k!}
= \frac{-2k!}{(k-1)! B_k} = \frac{-2k}{B_k}$。）

\begin{example}[$\sigma_{k-1}$ 的计算]
\label{ex:sigma-computation}
\leavevmode
\begin{itemize}
  \item $\sigma_3(1) = 1^3 = 1$，$\sigma_3(2) = 1^3 + 2^3 = 9$，
    $\sigma_3(3) = 1 + 27 = 28$，$\sigma_3(4) = 1 + 8 + 64 = 73$，
    $\sigma_3(6) = 1 + 8 + 27 + 216 = 252$。
  \item $\sigma_5(1) = 1$，$\sigma_5(2) = 1 + 32 = 33$，
    $\sigma_5(3) = 1 + 243 = 244$，$\sigma_5(4) = 1 + 32 + 1024 = 1057$。
\end{itemize}

验证 $E_4$ 的前几项：
$1 + 240(q + 9q^2 + 28q^3 + 73q^4 + \cdots)
= 1 + 240q + 2160q^2 + 6720q^3 + 17520q^4 + \cdots$ \quad $(\checkmark)$
\end{example}

\textbf{乘积法则与卷积。}
既然模形式相乘时权会相加，$E_rE_s$ 自然落在 $\Mk_{r+s}(\SL_2(\Z))$ 中。把两边都写成 $q$-展开，就会得到一串显式的卷积恒等式；这正是 CSS 第 I 部分特别强调的“把抽象结构落实到系数计算”的做法。

\begin{proposition}[Eisenstein 级数的卷积公式]
\label{prop:eisenstein-convolution}
设 $r,s\ge 4$ 为偶整数。则
\begin{align}
  E_r(\tau)E_s(\tau)
  &= 1
     - \frac{2r}{B_r}\sum_{n\ge 1}\sigma_{r-1}(n)q^n
     - \frac{2s}{B_s}\sum_{n\ge 1}\sigma_{s-1}(n)q^n \\
  &\qquad
     + \frac{4rs}{B_rB_s}
       \sum_{n\ge 1}\left(\sum_{m=1}^{n-1}
       \sigma_{r-1}(m)\sigma_{s-1}(n-m)\right)q^n.
\end{align}
因此 $E_rE_s-E_{r+s}$ 是权 $r+s$ 的尖点形式；特别地，若 $\Sk_{r+s}(\SL_2(\Z))=0$，则必有
\[
  E_rE_s=E_{r+s}.
\]
\end{proposition}

\begin{proof}
把 \cref{thm:Gk-q-expansion} 后面的标准化公式分别代入并直接相乘即可。由于 $E_rE_s$ 与 $E_{r+s}$ 的常数项都等于 $1$，所以差的常数项为 $0$，从而差是尖点形式。
\end{proof}

\begin{example}[从 $q$-展开验证乘积关系]
\label{ex:eisenstein-product-identities}
由
\[
  E_4=1+240q+2160q^2+6720q^3+\cdots
\]
直接平方，得到
\[
  E_4^2=1+480q+61920q^2+1050240q^3+\cdots.
\]
另一方面
\[
  E_8=1+480q+61920q^2+1050240q^3+\cdots,
\]
所以
\[
  E_4^2=E_8.
\]
同理，因为 $\Sk_{10}(\SL_2(\Z))=0$，也有 $E_4E_6=E_{10}$。
\end{example}

\begin{example}[权 $12$ 的第一条非平凡关系]
\label{ex:weight12-relation}
这时尖点空间不再消失，而是一维。把前几项乘出来：
\[
  E_4^3=1+720q+179280q^2+16954560q^3+\cdots,
\]
\[
  E_6^2=1-1008q+220752q^2+16519104q^3+\cdots.
\]
两式相减得到
\[
  E_4^3-E_6^2=1728q-41472q^2+435456q^3+\cdots.
\]
而
\[
  1728\Delta=1728q-41472q^2+435456q^3+\cdots,
\]
故
\[
  E_4^3-E_6^2=1728\Delta.
\]
这说明：从权 $12$ 开始，Eisenstein 部分与尖点部分真正开始发生相互作用。
\end{example}

\textbf{注。} 半整数权理论里也有与之平行的 Eisenstein 级数，最著名的是 Cohen--Eisenstein 级数 $H_{r+1/2}$。它们的 Fourier 系数不再是简单的除数和，而会涉及二次型表示数与虚二次域类数；从整数权的 $\sigma_{k-1}(n)$ 到半整数权的类数系数，正是很多后续算术现象的源头。

\begin{proposition}[$E_4, E_6$ 与 $\Delta$]
\label{prop:E4-E6-Delta}
\index{判别形式 (discriminant form)}\index{Delta@$\Delta(\tau)$}
\index{Ramanujan tau 函数}\index{tau@$\tau(n)$}
定义
\[
  \Delta(\tau) = \frac{E_4(\tau)^3 - E_6(\tau)^2}{1728}.
\]
则 $\Delta$ 是权 $12$ 的\textbf{尖点形式}（$\Delta(\infty) = 0$），
且有乘积公式
\[
  \Delta(\tau) = q \prod_{n=1}^{\infty}(1-q^n)^{24}
  = \sum_{n=1}^{\infty} \tau(n)\, q^n,
\]
其中 $\tau(n)$ 是 \textbf{Ramanujan $\tau$-函数}，$\tau(1) = 1$，$\tau(2) = -24$，
$\tau(3) = 252$，$\tau(4) = -1472$，$\ldots$
\end{proposition}

\begin{proof}
\textbf{$\Delta$ 是尖点形式}：$E_4(\infty) = 1$，$E_6(\infty) = 1$，故
$\Delta(\infty) = (1 - 1)/1728 = 0$，满足尖点条件。
权 $12 = k$ 由 $E_4^3$（权 $4 \times 3 = 12$）和 $E_6^2$（权 $6 \times 2 = 12$）的差给出。

\textbf{$\Delta$ 在 $\HH$ 上无零点（乘积公式）}：这比较深刻，
用到模形式空间的维数理论（见\cref{ch:dimension}）：
$\Sk_{12}(\SL_2(\Z))$ 恰好是一维的，
且 $\eta(\tau)^{24}$ 是其中一个非零元（Dedekind $\eta$ 函数
$\eta(\tau) = q^{1/24}\prod_{n=1}^\infty(1-q^n)$ 的 $24$ 次幂），
由比较常数项得 $\Delta(\tau) = q\prod_{n=1}^\infty(1-q^n)^{24}$。
\end{proof}

\textbf{$\Delta$ 与 Ramanujan 猜想。}
Ramanujan 在 1916 年观察到 $\tau(n)$ 的三个猜想：
\begin{enumerate}
  \item \textbf{乘性}：$\tau(mn) = \tau(m)\tau(n)$ 当 $\gcd(m,n) = 1$；
  \item \textbf{Hecke 关系}：$\tau(p^{r+1}) = \tau(p)\tau(p^r) - p^{11}\tau(p^{r-1})$；
  \item \textbf{估计}：$|\tau(p)| \leq 2p^{11/2}$。
\end{enumerate}
前两条由 Mordell（1917）证明，第三条（Ramanujan 猜想）由 Deligne 在 1974 年
利用 Weil 猜想的证明给出。乘性与 Hecke 算子有关（第三章），
估计与代数几何有关，这两个方向最终都汇入费马大定理的证明链。

\begin{example}[验证 $\tau$ 的乘性]
\label{ex:tau-multiplicative}
验证 $\tau(2)\tau(3) = \tau(6)$：
$(-24)(252) = -6048$，而 $\tau(6) = -6048$（$\checkmark$，需要计算 $q^6$ 系数）。

验证 $\tau(2)\tau(3) = \tau(6)$：
$(-24)(252) = -6048$，而 $\tau(6) = -6048$（$\checkmark$，需要计算 $q^6$ 系数）。

详细：$\Delta = q - 24q^2 + 252q^3 - 1472q^4 + 4830q^5 - 6048q^6 + \cdots$
\end{example}


% ============================================================================
